\section{Three purely quadratic labels in the actual arithmetic window}
\label{sec:pure-quadratic-actual-window}

\paragraph{Status and exact dependency.}
This conditional-interface proof uses the original-coefficient
strict theorem AC-G160 for the three labels with trivial cubic part.
It retains every actual omission in AC-G157-window and uses one
full local comparison cost for the enlarged union.
The separate class and local proofs are inputs at their recorded
scopes.  The remaining all-prime numerical estimate and mixed
realization are not supplied here.

Put $p=269$, $K=\mathbb Q(\omega)$, $t^3=p$ with $t$ real,
and let $\chi:G_K\to\mu_3$ describe $K(t)/K$.
Let $\phi$ be the primitive odd quartic character modulo sixteen
with $\phi(13)=i$ and kernel $\{1,7\}$.
Write $\delta_d=\chi_d|_{G_K}$ for $d\in\{-4,8,-8\}$.
For a fixed source datum, $0<\varepsilon<1/256$, an allowed
page and its prescribed pair $q,d$, keep
\[
 \begin{aligned}
 F_u&=\mathbb Q(\zeta_{2^{h(u)+2}},\zeta_7,
                          \sqrt{-q},\sqrt{-d},\omega),\\
 H_u&=F_u(t),\quad G_u=\operatorname{Gal}(H_u/K),\quad D_u=|G_u|,\\
 u^{2+8\varepsilon}&\leq2^{h(u)}\leq2u^{2+8\varepsilon},
                         \qquad T=\lfloor u^{399}\rfloor.
 \end{aligned}
 \tag{Q0W.1}
\]
The letter $d$ in the prescribed pair is independent of the
subscript on a quadratic character $\delta_d$.
All original $O_\ell,E_\ell$, normalizing ranks $b_\ell$,
full local conditions and derived indices are exactly
(FC.1)--(FC.2), (EU.2) and (TC.1).
The coefficient-prime range remains unrestricted.

For $u\geq2$ define the following actual disjoint sets:
\[
 \begin{aligned}
 Z_{16}^+&=\{(p,\chi^{-j}\phi^{-\sigma},e):j=1,2,\ \sigma=\pm1,\\
 &\hspace{2em}1\leq e<T,\ e\equiv-1\pmod{67}\},\\
 Z_4^+&=\{(p,\chi^{-j}\delta_{-4},e):j=1,2,
               \ 1\leq e<T,\ e\equiv-1\pmod{67}\},\\
 Z_8&=\{(p,\chi^{-j}\delta_d,e):j=1,2,\ d=8,-8,\\
 &\hspace{2em}1\leq e<T,\ e\equiv-1\pmod{134}\},\\
 Z_{\rm cub}^{(4)}&=\{(p,\chi^j,e):j=0,1,2,\ 1\leq e<T,\\
 &\hspace{2em}e\equiv0\pmod{268}\ \text{or}\ e\equiv-1\pmod{67}\},\\
 Z_3&=\{(3,1,e):1\leq e<T\},\\
 Z_{\rm quad}^0&=\{(p,\delta_d,e):d=-4,8,-8,\ 1\leq e<T,\\
 &\hspace{2em}e\equiv-1\pmod{134}\},\\
 Z^{\rm quad0}&=Z_{16}^+\sqcup Z_4^+\sqcup Z_8
                 \sqcup Z_{\rm cub}^{(4)}\sqcup Z_3\sqcup Z_{\rm quad}^0.
 \end{aligned}
 \tag{Q0W.2}
\]
Thus only $Z_{\rm quad}^0$ is added to the union of AC-G157.
Let $I_{399}^{\rm estquad0}$ be (TC.1) with precisely this new
union omitted; the old remaining sum is $I_{399}^{\rm estcub4}$.

\begin{theorem}[AC-G161-window: the three pure quadratic label reduction]
\label{AC-G161-window}
All labels in (Q0W.2) occur on the prescribed host, and their
strict groups on the original coefficient modules vanish.
For positive integer $T$ their exact disjoint counts are
\[
 \begin{aligned}
 |Z_{16}^+|&=4\lfloor T/67\rfloor,&|Z_4^+|&=2\lfloor T/67\rfloor,\\
 |Z_8|&=4\lfloor T/134\rfloor,&|Z_3|&=T-1,\\
 |Z_{\rm cub}^{(4)}|&=3\bigl(\lfloor(T-1)/268\rfloor
                                     +\lfloor T/67\rfloor\bigr),\\
 |Z_{\rm quad}^0|&=3\lfloor T/134\rfloor.
 \end{aligned}
 \tag{Q0W.3}
\]
Let $C_u$ be the actual normalized positive-index cost over the
entire union $Z^{\rm quad0}$.  Then
\[
 \begin{aligned}
 I_{399}&=I_{399}^{\rm estquad0}+C_u,\qquad
 I_{399}^{\rm estquad0}\leq I_{399}^{\rm estcub4}\leq I_{399},\\
 0\leq C_u&\leq8D_u^2T^2\log(8TD_u)
              +\tfrac12D_u\mathcal R_uT^2+2D_ur_uc_0T\\
 &=O_\varepsilon\bigl(u^{802+16\varepsilon}(\log u)^2\bigr).
 \end{aligned}
 \tag{Q0W.4}
\]
The first line is valid for nonnegative extended sums.
One full (FL.4) bound covers every old and new omitted label;
individual positive indices are not asserted zero.

The same-window premise (TC.3), or its AC-G157 version, may use
$I_{399}^{\rm estquad0}$ with all its simultaneous field,
coefficient, finite-character and weight quantifiers unchanged.
Either old premise implies the new one with unchanged parameters.
The new premise implies the full old one on the same $W,J$,
datum, page and pair, after only raising its initial threshold
and halving its positive relative slack.
The actual $R.t$ and every clause of the retained mixed input
$\mathcal R_{\rm cub}$ remain required and unproved.
\end{theorem}

\begin{proof}
\proofstep{1. Actual characters and the literal original modules.}
The conductor-eight cyclotomic field is contained in $F_u$
for $u\geq2$, so the three characters $\delta_d$ factor through
$G_u$.  They remain distinct and nontrivial after restriction to
$K$, because $K\cap\mathbb Q(\zeta_8)=\mathbb Q$.
Also $F_u\cap K(t)=K$: an abelian extension of $\mathbb Q$
cannot contain the nonabelian $S_3$ splitting field $K(t)$,
and $K(t)/K$ has prime degree three.
Consequently the actual cubic subgroup of $G_u$ distinguishes
each pure quadratic character from every listed character with
nontrivial cubic factor.  The remaining pure-cubic label $1$
is distinct because each $\delta_d$ is nontrivial.
The coefficient-three family has a different coefficient prime.
This proves the new disjointness; all old occurrence and
disjointness statements of AC-G157 remain unchanged.

The original $O_p$ contains the values $\pm1$.
For each new label its coefficient is literally
\[
 \begin{aligned}
 A_{d,e}&=(O_p\otimes_{\mathbb Z_p}\mathbb Q_p/\mathbb Z_p)
                                                    (\delta_d)(-e)\\
        &=(E_p/O_p)(\delta_dc_p^{-e}).
 \end{aligned}
 \tag{Q0W.5}
\]
Apply (PQ.2) of AC-G160 with $m=0$ and $\kappa=1$.
It permits this actual original $O_p$ with arbitrary ramification
and proves strict zero for every $e\equiv133,267\pmod{268}$,
exactly $e\equiv-1\pmod{134}$.
That source theorem proves the trivial cubic part directly.
It is not an extension of a theorem restricted to $j=1,2$.
Its quadratic character at $p$ occurs only in the residual
cover for the first progression; it is absent from (Q0W.5).
Every old strict zero, including the coefficient-three statement
without averaging by three, is retained from AC-G157.

\proofstep{2. Full local conditions and counts.}
The strict theorem imposes zero in the full group
$H^1(G_{K_v},A_{d,e})$ at each $v\mid p$, and zero in
$H^1(I_v,A_{d,e})$ at every other finite place.
These are the original (FC.2) local conditions.
The coefficient (Q0W.5) is unramified away from $2,p$.
At any added away-$p$ prime of $\Sigma_p$, inertia acts trivially;
a coboundary restriction there is zero.
The cocycle identity then makes it zero on every conjugate inertia
subgroup and their generated closed normal subgroup.
Both the action and cocycle descend to the smaller ramification
set.  Conversely inflation is injective and preserves all the
remaining local conditions.  No full condition at $p$ is weakened.
Thus the global theorem applies to exactly the strict modules
used by (FC.2), also on the enlarged ramification set of the host.

The new weights have the form $e=134n-1$, $n\geq1$.
The inequality $1\leq e<T$ is equivalent to $1\leq n\leq T/134$.
For each of three distinct finite labels there are
$\lfloor T/134\rfloor$ weights.
This proves the last count of (Q0W.3); every other count is the
unchanged AC-G157 count.  Precisely these added strict terms
vanish in (FC.15).  This does not discard their derived-index costs.

\proofstep{3. One total-union comparison cost.}
Use the original dual complex
$C_{\ell,\eta,e}=R\Gamma(G_{K,\Sigma_\ell},
O_\ell(\eta^{-1})(e+1))$, with the codifferent normalization
of (FL.2).  For every removed label, (EU.2) and strict zero give
\[
 \begin{aligned}
 \frac{\log\#k_{O_\ell}}{b_\ell}\max(0,i_{\ell,\eta,e})
 &\leq\frac{\log\#H^2(C_{\ell,\eta,e})}{b_\ell}\\
 &=\frac{|\log\#H^2(C_{\ell,\eta,e})
                  -\log\#\mathcal S_{\ell,\eta}(e)|}{b_\ell}.
 \end{aligned}
 \tag{Q0W.6}
\]
The positive index and the entire original normalizing rank
$b_\ell$ are retained.  At each weight the whole union is one
subsum of the nonnegative all-prime absolute error (FL.4),
so its cost is at most
\[
                  16eD_u^2\log(8eD_u)+D_ue\mathcal R_u+2D_ur_uc_0.
 \tag{Q0W.7}
\]
Summing once over $1\leq e<T$, using $\sum e\leq T^2/2$
and $\log(8eD_u)\leq\log(8TD_u)$, proves (Q0W.4).
No extra full error is charged for the new family.
Nonnegative additivity proves its extended-sum identity and
monotonicity without subtracting infinite quantities.

The unchanged host estimates are $D_u\leq288u^{2+8\varepsilon}$,
$r_u\leq12$ and
$\mathcal R_u=O_\varepsilon(u^{\varepsilon/5}(\log u)^2)$.
Together with $T\leq u^{399}$ they give largest exponents
$802+16\varepsilon$ and $800+41\varepsilon/5$.
Hence $C_u=o(\exp(u^\varepsilon/2))$, uniformly over the
prescribed page and pair, for each fixed positive $\varepsilon$.

\proofstep{4. The original common window and slack.}
Fix the retained datum, page and $0<\varepsilon<1/256$.
Suppose the new premise supplies $J\geq1$, $X_0$ and
$s\in(0,1)$ independent of $X,q,d,u,\ell,\eta,e$, such that
for every $X\geq X_0$ and prescribed
$\operatorname{WinPair}(\varepsilon,X^{1/J},q,d)$ there is
one $W$ with $X^{1/J}\leq W$, $W^2\leq X$ and
\[
 I_{399}^{\rm estquad0}(H_u,u)\leq(1-s)\exp(u^\varepsilon/2)
                                      \quad(W\leq u\leq W^2).
 \tag{Q0W.8}
\]
Choose $U_0\geq2$ making (Q0W.4) at most
$(s/2)\exp(u^\varepsilon/2)$ for every $u\geq U_0$.
Only the retained datum, page, $\varepsilon$ and slack enter
this uniform choice.  Replace $X_0$ by $\max(X_0,U_0^J)$.
The very same $W$ then gives
\[
 I_{399}(H_u,u)\leq(1-s/2)\exp(u^\varepsilon/2)
                                      \quad(W\leq u\leq W^2).
 \tag{Q0W.9}
\]
For a fixed finite simultaneous family use the maximum of its
uniform error constants while keeping this one $W$.
There is no intersection of separately chosen windows.
Conversely either old premise implies (Q0W.8) by nonnegativity,
with unchanged parameters.  Uniform little-$o$ versions agree.
Every host, original coefficient, actual source $R.t$, finite
comparison cost and mixed realization clause remains as before.
\end{proof}

\paragraph{Exact remaining gate.}
The new set is a proper additional subset of actual terms for
$T\geq134$; no strictly smaller numerical value is asserted.
Every unlisted prime, finite character and weight remains in the
unknown sum.  The all-prime numerical bound, mixed theta/Selmer
identification and analytic realization, and H04.4 remain open.
Only the previously audited SGA~1, Expos\'e~V, \S\S4 and 7
finite-cover and original rank-one Artin--Tate dictionary is used.
No new motive, comparison theorem, maximal vanishing classification
or Lean execution is asserted.
